(1) 颈部空气向外移动 x，空腔体积增加 \(Ax\)，绝热微扰给 \(\delta p=-\rho c^2Ax/V\)。作用在截面积 A 上的净回复力
\[
F=-\frac{\rho c^2A^2}{V}x=-kx,\qquad k=\frac{\rho c^2A^2}{V}.
\]
腔体气体的可压缩性充当弹簧。

(2) 颈部有效质量 \(m_{\rm eff}=\rho A L_e\)，故
\[
\omega_H=\sqrt{\frac{k}{m_{\rm eff}}}
=c\sqrt{\frac{A}{VL_e}},\qquad f_H=\frac{\omega_H}{2\pi}.
\]
这个频率来自颈部气柱惯性与空腔声压回复力，量纲为逆时间。

(3) 加线性黏滞阻力后，动力学方程
\(m_{\rm eff}\ddot x+b\dot x+kx=0\)。
记 \(\gamma=b/(2m_{\rm eff})\)，当 \(\gamma<\omega_H\) 时，欠阻尼角频率
\[
\omega_d=\sqrt{\omega_H^2-\gamma^2}.
\]
临界阻尼系数 \(b_c=2\sqrt{km_{\rm eff}}\)。机械能 \(E=m_{\rm eff}\dot x^2/2+kx^2/2\) 满足 \(\dot E=-b\dot x^2\le0\)。

(4) 外界声压谐波幅值 pext 对颈部产生的驱动力幅值 \(F_0=A p_{\rm ext}\)。采用 \(e^{-i\omega t}\) 约定，稳态位移复幅
\[
X=\frac{A p_{\rm ext}}{k-m_{\rm eff}\omega^2-i b\omega}.
\]
平均阻尼耗散功率 \(\overline P=\frac12b\omega^2|X|^2\ge0\)。频率接近 \(\omega_H\) 且阻尼较小时响应显著增强；假设只在波长远大于腔体尺度的集总声学范围成立。
